\documentclass[11pt]{article}
\usepackage[margin=1.05in]{geometry}
\usepackage{amsmath,amssymb,amsthm,mathtools,booktabs,array,microtype}
\usepackage[hidelinks]{hyperref}
\newtheorem{theorem}{Theorem}[section]
\newtheorem{lemma}[theorem]{Lemma}
\newtheorem{corollary}[theorem]{Corollary}
\theoremstyle{remark}\newtheorem{remark}[theorem]{Remark}
\newcommand{\ord}{\operatorname{ord}}
\newcommand{\Rea}{\operatorname{Re}}
\newcommand{\Frob}{\operatorname{Frob}}
\newcommand{\Q}{\mathbb Q}
\newcommand{\C}{\mathbb C}
\newcommand{\sgn}{\operatorname{sgn}}
\newcommand{\eps}{\varepsilon}
\title{Reflection detectors for nonabelian Riemann hypotheses\\\large A one-sided prime-sum transfer from the Alp\"oge--Claude proof}
\author{Research note for ProveIt}
\date{25 November 2026}
\begin{document}
\maketitle
\begin{abstract}
The compact Lean proof of Dirichlet GRH separates a geometric height inequality from an analytic one-sided criterion. We combine the latter with local positivity of Heilbronn characters and reflection factorizations. For an $S_n$-Galois extension $L/\Q$, let $E=L^{\langle(12)\rangle}$. At every point $\rho$ in the open right half of the critical strip, assuming Dirichlet GRH,
\[
\ord_\rho\zeta_L\le 2(n-H_n)\ord_\rho\zeta_E,
\qquad H_n=\sum_{k=1}^n\frac1k.
\]
No Artin holomorphy assumption is made. Thus a single one-sided, centered, smoothed rational-prime test involving only the identity and transposition Frobenius classes is equivalent to RH for $\zeta_L$, given Dirichlet GRH. An analogous result for finite real reflection groups uses one representative per reflection class, combined into a single scalar test. We give proofs, a reproducible $S_5$ example, and the exact analytic interface that would need to be supplied by an extension of the geometric argument. Neither the new prime inequality nor full nonabelian GRH is proved here; priority of the multiplicity bound is not established.
\end{abstract}
\tableofcontents

\section{Source audit and scope}
The source repository \cite{repo} describes the work as a shorter Lean rewriting of Levent Alp\"oge's proof with Claude. It reports kernel and Comparator checks and the three standard Lean axioms, but also explicitly says its text had not been read through by a mathematician. The audit here is static, through the GitHub connector; the reported checks were not independently rerun. The inspected tree is \texttt{72fead102f663d5c7a407964a13dc0cf9d977207}. Blob identifiers and links are in \texttt{source-audit.txt}. The root README pins Lean 4.35.0 and Mathlib at \texttt{56a005855ea17e0b0958f73ec801ae140b0a1679}.

The five statements in \texttt{grh-compact/src/Main.lean} are RH, its zeta-zero formulation, nonvanishing for Dirichlet $L$-functions right of $1/2$ with the principal pole excepted, their zero-line formulation in the critical strip, and the primitive special case. The final theorem is not RH for all number fields or all automorphic $L$-functions. We use Dirichlet GRH as an established input in the deductions below, following the requested premise.

\paragraph{Architecture.} The guide and \texttt{New/Main.lean} identify three ingredients: short forms, enough forms, and lower bounds for integral forms. A common covolume and its index are compared above and below; the main height cancels, yielding \texttt{star\_inequality}, namely $m\,\mathrm{remainder}\le (\log b)^2\,\mathrm{budget}$. Here the number $m$ of lifts is of order $\sqrt{x}$. The second-scale argument in \texttt{OneSidedCriterion.lean} supplies an appropriate long window, and the lemma\newline \texttt{eventually\_primeError\_le\_rpow} gives a one-sided bound with exponent arbitrarily close to $-1/2$. The final Landau criterion uses quadratic twists to avoid cancellation. The relevant kernel is $\left(1/\nu-1/x\right)_+$, with coefficients $\Lambda(\nu)(1-\chi_d(\nu))$ restricted to $1\bmod M$. The exact identity appears in \texttt{PrimeError.lean}.

The file \texttt{PolarLandau.lean} implements a general one-sided analytic implication for functions \texttt{IsPolar}, with nonnegative real residues. This is the most directly reusable analytic theorem. The class-field file \texttt{FrobeniusFilter.lean} supplies existence of degree-one Frobenius primes, not a quantitative Chebotarev error. \texttt{HeckeContinuation.lean} proves a restricted continuation near $1$, not a general Hecke GRH. These distinctions matter when defining an extension.

\paragraph{Earlier reports.} The incoming directory \cite{incoming} contains prior work on CM and sparse singular moduli, class-group exponents, cubic and quintic fields, Sylow transfers, $A_5$ detectors, and sampling and ramification. We inspected the directory and central-directory filenames of several smaller ZIPs; see the audit. We did not read all manuscript contents, and could not establish nonoverlap or historical priority. To avoid an obvious repetition, the new focus is the symmetric and real-reflection-group local character argument, with a quantitative multiplicity comparison. This is not a list of conventional applications of Dirichlet GRH.

\section{A multiplicity detector without Artin holomorphy}
Let $L/\Q$ be finite Galois, $G=\operatorname{Gal}(L/\Q)$, and fix $\rho$ with $1/2<\Rea\rho<1$. Write $m=\ord_\rho\zeta_L\ge0$. By Brauer induction Artin $L$-functions are meromorphic, so their orders $n_\chi$ are integers (negative orders are allowed). Define the Heilbronn virtual character
\[
 \Theta=\sum_{\chi\in\operatorname{Irr}(G)} n_\chi\chi,
 \qquad n_\chi=\ord_\rho L(s,\chi,L/\Q).
\]
Artin formalism, Frobenius reciprocity and the regular character give, for every subgroup $U$,
\begin{equation}\label{eq:orders}
 \Theta(1)=m,\quad
 \langle\Theta,\operatorname{Ind}_U^G1\rangle
 =\langle\Theta|_U,1\rangle=\ord_\rho\zeta_{L^U}.
\end{equation}
These are classical identities (see \cite{heilbronn,foote}). Furthermore, for a one-dimensional character $\eta$ of $U$, $\langle\Theta|_U,\eta\rangle$ is the order of a finite-order Hecke $L$-function over $L^U$ at $\rho$, hence nonnegative. If every irreducible character of $U$ is induced from a linear character of a subgroup, the same statement holds for every irreducible of $U$. Thus the restriction is an actual character of dimension $m$ on every abelian or dihedral subgroup. For dihedral groups this follows directly by inducing the characters of the cyclic rotation subgroup in the non-inversion-invariant cases; the other irreducibles are one-dimensional. We do not assert that $\Theta$ itself is an actual character.

For any linear $\epsilon:G\to\{\pm1\}$, its Artin $L$-function over $\Q$ is a Dirichlet $L$-function (including the trivial case). Dirichlet GRH implies
\begin{equation}\label{eq:signorth}
 \langle\Theta,\epsilon\rangle=0.
\end{equation}

\begin{theorem}[Local reflection inequality]\label{thm:general}
Suppose $T$ is a conjugation-invariant set of involutions in $G$ and $\epsilon$ is a linear character with $\epsilon(t)=-1$ for $t\in T$. Assume each $g$ has an expression $g=uv$, where $u$ and $v$ are products of, respectively, $r_g$ and $s_g$ pairwise commuting members of $T$ (empty products are allowed). Put
\[
 a_t=\ord_\rho\zeta_{L^{\langle t\rangle}}\quad(t\in T).
\]
For a chosen expression of each $g$ with reflection list $t_{g,1},\dots,t_{g,r_g+s_g}$, one has
\begin{equation}\label{eq:weighted}
 m\le \frac{2}{|G|}\sum_{g\in G}\sum_{j=1}^{r_g+s_g}a_{t_{g,j}}.
\end{equation}
In particular if every expression has at most $d$ reflections, then $m\le2d\max_{t\in T}a_t$. If all $a_t=0$, then $m=0$.
\end{theorem}
\begin{proof}
On an abelian subgroup containing $t$, the actual representation affording $\Theta$ has $+1$ eigenspace of dimension $a_t=(m+\Theta(t))/2$, by \eqref{eq:orders}. For $u=t_1\cdots t_r$ with commuting factors, the intersection of their $-1$ eigenspaces has codimension at most $\sum a_{t_i}$. Hence the dimension $b(u)$ of the eigenspace on which $u$ differs from the scalar $\epsilon(u)$ is at most $\sum a_{t_i}$. This number equals $(m-\epsilon(u)\Theta(u))/2$ and depends only on the character value, not on the chosen subgroup representation. The same applies to $v$.

Now $\langle u,v\rangle$ is dihedral, cyclic of order at most two, or trivial. Indeed $u,v$ square to one, $uv$ generates a cyclic normal subgroup and conjugation by $u$ inverts $uv$. Use a unitary representation of this subgroup affording $\Theta$. After twisting by $\epsilon$, the intersection of the fixed spaces of $u$ and $v$ has codimension at most $b(u)+b(v)$. Their product fixes this intersection. All other eigenvalues of the unitary product have real part at least $-1$. Consequently
\[
 m-\Rea(\epsilon(g)\Theta(g))\le2(b(u)+b(v))
 \le2\sum_j a_{t_{g,j}}.
\]
Average over $G$. The average of $\epsilon(g)\Theta(g)$ is zero by \eqref{eq:signorth}, proving the claim. The argument only compares characters on actual local representations, never assumes simultaneous global positive definiteness.
\end{proof}

\begin{corollary}[Symmetric groups]\label{cor:sn}
For $G=S_n$ ($n\ge2$), $E=L^{\langle(12)\rangle}$ and $H_n=1+1/2+\cdots+1/n$,
\[
 \ord_\rho\zeta_L\le 2(n-H_n)\ord_\rho\zeta_E.
\]
In particular RH for $\zeta_E$ implies RH for $\zeta_L$ under Dirichlet GRH. The reverse implication follows from the factorization for the quadratic extension $L/E$ (or Aramata--Brauer).
\end{corollary}
\begin{proof}
For an $h$-cycle index the elements by $\mathbb Z/h\mathbb Z$ with the cycle acting by $i\mapsto i+1$. Let $u(i)=1-i$, $v(i)=-i$. Then $uv(i)=i+1$. For odd $h$, both involutions have $(h-1)/2$ disjoint transpositions; for even $h$, their numbers are $h/2$ and $h/2-1$. On distinct cycles the factors commute. Thus $r_g+s_g=n-c(g)$, where $c(g)$ counts cycles including fixed points. All transpositions are conjugate; use $\epsilon=\sgn$. A uniform random permutation on $n$ letters has expected $H_n$ cycles: when inserting the last letter into the cycle notation of a permutation of $n-1$ letters, one of $n$ equally likely positions creates a new singleton and the others insert it after an existing letter. Apply \eqref{eq:weighted}. The quadratic factor is a Hecke $L$-function holomorphic in the strip, so $\zeta_E$ zeros persist in $\zeta_L$.
\end{proof}

\begin{remark}[How strong is the constant?]
For $n\ge3$, the genuine character $\Theta=\sgn\otimes\mathrm{Std}$, with $\mathrm{Std}$ the standard $(n-1)$-dimensional representation, has $m=n-1$, $a_t=1$, and $\langle\Theta,\sgn\rangle=0$ (also $\langle\Theta,1\rangle=0$). Therefore no inequality $m\le C_n a_t$ derived solely from local positivity and the vanishing of the two linear-character coefficients can have $C_n<n-1$. This is a character example, not a claim that such orders occur at an actual zero. Our upper coefficient $2(n-H_n)$ is within a factor two of this obstruction. For $n=5$, the coefficient is $163/30\approx5.4333$, so a zero of order $m$ in $\zeta_L$ must have order at least $\lceil30m/163\rceil$ in $\zeta_E$.
\end{remark}

\begin{corollary}[Finite real reflection groups]\label{cor:cox}
Let $G=W$ be a finite real reflection group of rank $d$, acting in its faithful real reflection representation. Choose one reflection $t_j$ from each conjugacy class of reflections, and write $E_j=L^{\langle t_j\rangle}$, $1\le j\le q$. Then
\[
 \ord_\rho\zeta_L\le 2d\max_j\ord_\rho\zeta_{E_j}
 \le 2d\sum_j\ord_\rho\zeta_{E_j}.
\]
RH for all $\zeta_{E_j}$ is equivalent to RH for $\zeta_L$, given Dirichlet GRH. There is one detector for simply laced irreducible Weyl groups, two for types $B_n$ ($n\ge2$), $F_4$, and $I_2(2k)$ ($k\ge2$).
\end{corollary}
\begin{proof}
The classical finite Coxeter factorization results of Carter and Springer \cite{carter,springer,hurwitz} give $w=uv$ with each of $u,v$ an involution or identity, each a product of commuting reflections, and with total reflection length $\ell_T(w)=\operatorname{codim}V^w\le d$. We use the length-additive involution factorization theorem in the form recorded in \cite{hurwitz}. The determinant is $-1$ on every reflection. Apply Theorem~\ref{thm:general}. Each $L/E_j$ is quadratic and its Hecke factor is holomorphic in the strip, giving the reverse implication. The reflection-class counts follow from conjugacy of simple reflections along odd-labelled paths in the Coxeter diagram.
\end{proof}
\paragraph{Conservative variant.} Even without the length-additive theorem, the product-of-two-involutions theorem and commuting-reflection decomposition give at most $2d$ factors and hence the same zero-detection conclusion with constant $4d$. The $S_n$ result needs neither classification nor finite computation.

\section{The one-sided prime criterion}
For each fixed number field $E$ let
\[
 \psi_E(x)=\sum_{(\mathfrak p,k):\,N\mathfrak p^k\le x}\log N\mathfrak p,
 \quad A_E(x)=\psi_E(x)-x,
 \quad \kappa_E=\int_1^\infty A_E(u)\frac{du}{u^2}.
\]
The fixed-field prime ideal theorem with a classical zero-free region gives absolute convergence of the integral (a possible fixed exceptional real zero $<1$ causes an integrable term). Define the centered smoothed remainder
\begin{align}
 R_E(x)&=\int_1^x A_E(u)\frac{du}{u^2}-\kappa_E\\
 &=\sum_{N\mathfrak p^k\le x}\log N\mathfrak p\left(\frac1{N\mathfrak p^k}-\frac1x\right)
      -\log x-\kappa_E.\label{eq:kernel}
\end{align}
This is precisely the same smoothing shape as the upstream \texttt{primeError}, although the arithmetic coefficients differ. From $B_E(s)=-\zeta_E'(s)/\zeta_E(s)$ and partial summation,
\begin{equation}\label{eq:trans}
 \int_1^\infty R_E(x)x^{-s}\,dx
 =\frac{B_E(s)/s-1/(s-1)-\kappa_E}{s-1}
 =:\Phi_E(s)\qquad(\Rea s>1).
\end{equation}
Indeed integration by parts uses $R_E(1)=-\kappa_E$ and $R'_E(x)=A_E(x)/x^2$ almost everywhere. The constant is $\kappa_E=-1-\gamma_E$, where $\gamma_E$ is the Euler--Kronecker constant defined by $\zeta_E(s)=c_E((s-1)^{-1}+\gamma_E+O(s-1))$: take the limit of $B_E(s)/s-1/(s-1)$ as $s\downarrow1$ using dominated convergence. Thus the singularity at $1$ of \eqref{eq:trans} is removable. The transform is globally meromorphic and every pole in the right half of the critical strip is simple with residue
\begin{equation}\label{eq:res}
 \operatorname{Res}_{s=\rho}\Phi_E(s)
 =\frac{\ord_\rho\zeta_E}{\rho(1-\rho)}.
\end{equation}
For a real $\rho\in(1/2,1)$ this residue is nonnegative. This sign is critical: reversing the inequality would allow an exceptional real zero to escape detection.

\begin{lemma}[Positive-transform argument]\label{lem:landau}
For a finite nonempty collection of number fields $E_j$ and positive real weights $w_j$, let $R=\sum_j w_jR_{E_j}$. If for some $1/2<\beta<1$ and sufficiently large $x$,
\[
 R(x)\le C x^{\beta-1},
\]
then no $\zeta_{E_j}$ has a zero in $\Rea s>\beta$ (excluding the pole at $1$).
\end{lemma}
\begin{proof}
Increase $C$ to make $f(x)=Cx^{\beta-1}-R(x)\ge0$ on all $x\ge1$. This is possible because the remainders are continuous on compact intervals. The Mellin transform $M(s)=\int_1^\infty f(x)x^{-s}dx$ initially agrees for $\Rea s>1$ with $C/(s-\beta)-\sum_j w_j\Phi_{E_j}(s)$, and hence continues meromorphically. If its convergence abscissa $\sigma_c$ exceeds $\beta$, it must be singular at the real point $\sigma_c$ by Landau's positivity lemma. For completeness, changing variables $x=e^v$ gives a Laplace integral of a nonnegative measure on $[0,\infty)$. If the transform is analytic at the abscissa, take a real point just to its right inside a disk of analyticity extending left of the abscissa. Expanding $e^{hv}$ with nonnegative coefficients and using monotone convergence on the Taylor series extends convergence to a real point left of the abscissa, a contradiction.

At $s=1$ and above the continuation has no pole. At a real pole $\sigma_c\in(\beta,1)$, \eqref{eq:res} implies that the residue of $M$ is strictly negative, which would force $M(\sigma)\to-\infty$ as $\sigma\downarrow\sigma_c$. But $M(\sigma)\ge0$ on its convergence interval. Thus $\sigma_c\le\beta$ (the case of abscissa $-\infty$ is harmless). The integral is holomorphic in $\Rea s>\beta$ by uniform convergence on compacta. A zero of any $\zeta_{E_j}$ at $\rho$ in this region would yield residue $-\sum_j w_j\ord_\rho\zeta_{E_j}/(\rho(1-\rho))\ne0$, contradicting holomorphy. All zero orders are nonnegative and the weights are positive, so cancellation is impossible even at nonreal zeros.
\end{proof}

\section{A single rational-prime inequality for an entire Galois closure}
\label{sec:prime}
Let $T$ be the reflection classes of Corollary~\ref{cor:cox}, or just the transposition class in Corollary~\ref{cor:sn}. For $p$ unramified in $L$, $r_j(p)$ is the number of fixed cosets of $\Frob_p$ acting on $G/\langle t_j\rangle$. Equivalently it is the number of degree-one primes of $E_j$ above $p$. This is computable from the splitting of a defining polynomial for $L$ if the Galois group and conjugacy type can be determined. Put $\theta_j(x)=\sum_{p\le x,\,p\ \mathrm{unram.}} r_j(p)\log p$ and define $R_j^{(1)}$ by replacing $\psi_{E_j}$ in \eqref{eq:kernel} with $\theta_j$ and using the corresponding integral constant $\kappa_j^{(1)}$.

\begin{lemma}[Prime-power removal]\label{lem:primepower}
For fixed $L,E_j$, $R_j^{(1)}(x)-R_{E_j}(x)=O_{L,E_j}(x^{-1/2}\log x)$ as $x\to\infty$.
\end{lemma}
\begin{proof}
All unramified contributions omitted from $\theta_j$ involve $p^{fk}\le x$ with $fk\ge2$, hence $p\le\sqrt{x}$. For each $p$ the total such contribution is at most $[E_j:\Q]\log x$ using the degree identity $\sum_{\mathfrak p|p}e_{\mathfrak p}f_{\mathfrak p}=[E_j:\Q]$. The finitely many ramified $p$ give $O_L(\log x)$ as well. Thus the difference is $O(\sqrt{x}\log x)$. The centered remainders are negative tail integrals of their errors divided by $u^2$, and integrating this majorant gives the stated bound. In particular the two centering constants need not be equal; ignoring this difference would invalidate an eventual decaying test.
\end{proof}

\begin{theorem}[Reflection prime test]\label{thm:test}
Under Dirichlet GRH, for any Galois extension of $\Q$ with finite real reflection group $G$, the following are equivalent:
\begin{enumerate}
\item RH for $\zeta_L$;
\item for every $\eta>0$, a constant $C_{L,\eta}$ exists such that for all sufficiently large $x$,
\begin{equation}\label{eq:target}
 \sum_{j=1}^q w_j R_j^{(1)}(x)\le C_{L,\eta}x^{-1/2+\eta},
\end{equation}
where the fixed weights $w_j>0$ are arbitrary.
\end{enumerate}
Only the single summed upper bound is required, not separate two-sided estimates. For $G=S_n$, $q=1$, and the test uses exactly the weight
\begin{equation}\label{eq:weight}
 r(p)=\begin{cases}
 n!/2,&\Frob_p=1,\\
 (n-2)!,&\Frob_p\text{ a transposition},\\
 0,&\text{otherwise}.
 \end{cases}
\end{equation}
\end{theorem}
\begin{proof}
By Lemma~\ref{lem:primepower}, condition 2 also holds for the full prime-power remainders. If $\zeta_L$ has a zero with $\Rea\rho>1/2$ in the strip, choose $\eta$ strictly between zero and $\Rea\rho-1/2$. Corollary~\ref{cor:cox} implies some $E_j$ has a zero at the same point, contradicting Lemma~\ref{lem:landau} with $\beta=1/2+\eta$. The functional equation gives RH. Conversely, if RH holds for $L$, the quadratic factorization implies RH for every $E_j$. The classical explicit formula for a fixed Dedekind zeta function gives $\psi_{E_j}(x)=x+O_{E_j}(\sqrt{x}\log^2 x)$ under its RH; integration of the tail gives $R_{E_j}(x)=O_{E_j}(x^{-1/2}\log^2 x)$. Lemma~\ref{lem:primepower} now implies 2.

Finally, the fixed-coset count is $|\{a\in G:a^{-1}ga\in\langle t\rangle\}|/2$. At $g=1$ this is $|G|/2$, for $g$ conjugate to $t$ it is $|C_G(t)|/2$, and otherwise zero. For a transposition $|C_{S_n}(t)|=2(n-2)!$; this proves \eqref{eq:weight}. Its average over $G$ is one by Burnside's lemma. In particular subtracting $\log x$ rather than a density depending on $n$ is correct.
\end{proof}
\paragraph{Contrapositive.} A hypothetical zero of $\zeta_L$ of real part $\sigma>1/2$ forces the left side of \eqref{eq:target}, divided by $x^{\beta-1}$, to be unbounded above on every tail for every $1/2<\beta<\sigma$. This is a one-sided oscillation consequence of a hypothetical counterexample, not an unconditional asymptotic estimate.

\section{Example: the normal closure of a quintic}
Let $f(X)=X^5-X-1$, $L$ its splitting field. The discriminant is $2869=19\cdot151$; the factor degrees modulo $2$ are $(2,3)$ and modulo $3$ are $(5)$. These exact statements are verified by the accompanying Python program. Irreducibility modulo $3$ proves irreducibility over $\Q$. The unramified Frobenius classes supply a $5$-cycle and a permutation of type $(2,3)$, whose cube is a transposition. The conjugates of a transposition by a $5$-cycle generate $S_5$: its edges on the five letters form a connected graph, and transpositions on the edges of a connected graph generate its full symmetric group. Thus $G=S_5$. In particular $[E:\Q]=60$ and $[L:\Q]=120$.

For every $p\nmid2869$ let $d(p)$ denote the partition of $5$ consisting of factor degrees of $f$ modulo $p$. The sole weight is $60$ for $d(p)=(1,1,1,1,1)$, $6$ for $d(p)=(1,1,1,2)$, and zero otherwise. All other factor patterns may be discarded in the arithmetic sum. The density-weight check is $60/120+6\cdot10/120=1$. A single upper estimate on the centered smoothed sum of these weights suffices for RH of the degree-120 normal closure. The factor-degree observations are finite and do not certify this estimate. The test is not claimed easier computationally than all existing approaches; it specifies a concrete arithmetic interface for a possible height argument.

\section{What can actually be ported to Lean?}
The source \texttt{IsPolar} asks for measurability, a growth bound, a global meromorphic transform, analyticity of $(s-1)\Phi$ for $\Rea s\ge1$, and a simple-pole residue limit for $\Rea s>0$. For full prime-power sums $v\mapsto R(e^v)$, \eqref{eq:trans} gives the required transform exactly. It also gives the residue hypothesis of \texttt{IsPolar.landau} when $\beta>1/2$. This is a more literal reuse than a reformulation of the Dirichlet conclusion. The prime-only transform need only continue to $\Rea s>1/2$ because of the omitted higher powers; apply the existing global lemma before removal, not after it. Its eventual bound is written with unit coefficient in Lean; divide $R$ by a positive bound constant (or absorb the constant by increasing the exponent slightly). The fixed-field prime ideal theorem, Dedekind logarithmic derivative, analytic continuation, and order identities must be proved or imported with exact statements. For $0<\Rea s\le1/2$ the simple-pole residue limit is still valid away from $s=1$; $s=0$ is outside the hypothesis.

The group-theoretic proof has its own formalization route: (i) finite virtual characters, restriction, induction and orthogonality; (ii) positivity from finite-order Hecke holomorphy and monomiality of dihedral groups; (iii) codimension bounds for intersection of subspaces in unitary representations; (iv) permutation decompositions with explicit cycle indexing; (v) trace averaging. For arbitrary Coxeter groups import or prove the length-additive involution factorization theorem. The included \texttt{lean-interface.md} is a precise dependency plan, not a compiled Lean proof. No Lean binary was available here.

\section{Remaining geometric problem and research questions}
The source does not establish \eqref{eq:target}. Its arithmetic signal is a quadratic splitting character on rational primes in a residue class. Its CM local divisibility and theta-column arguments are specific to that arithmetic data. Our $r_j(p)$ are nonabelian Frobenius permutation weights. It is not justified to replace one by the other in \texttt{lambdaClass} and assume \texttt{star\_inequality} survives. In particular the coefficients in \eqref{eq:weight} are not the quadratic nonsplit indicators used upstream.

A viable next investigation would start with the $S_5$ example: design a global arithmetic object whose local divisibility gains are proportional to the two nonzero weights; prove an upper covolume and lower covolume estimate with the same main term; quantify missing dimension and index losses; then match the centered kernel and preserve the sign in \eqref{eq:target}. The upstream long-window lemma may be useful for the second scale once such a geometry exists. The square-root saving in the original proof is tied to the multiplicity of theta lifts; it cannot be imported without the corresponding estimates.

Other questions: Can the coefficient $2(n-H_n)$ be improved to $n-1$ using only the Heilbronn cone? Can one replace the degree-$n!/2$ subfield by lower-degree subfields without invoking Artin holomorphy? Which non-reflection Galois groups admit similar two-involution decompositions and a linear detector? Can local positivity on larger monomial subgroups produce an analogous criterion for groups such as $A_n$, where the determinant character trick is unavailable? These are genuinely separate group-theoretic questions from the geometric prime bound.

\paragraph{Claim ledger.} Proved here: the local-character inequality conditional on standard Artin--Hecke formalism and Dirichlet GRH; the symmetric and reflection group deductions; and the one-sided analytic equivalence. Verified computationally: finite permutation identities and the explicit quintic factorization data. Not proved: the geometric inequality, any previously unresolved case of full Dedekind GRH, Artin holomorphy, or formal Lean acceptance of the new work. The proof is a research derivation with classical inputs, not an independently peer-reviewed result; novelty is plausible as a combination but not certified.

\begin{thebibliography}{9}
\bibitem{repo} L. Sernau et al., \emph{GRH for Dirichlet L-functions, in Lean 4}, repository, 25 Nov. 2026. \url{https://github.com/lsernau/grh-lean}. See also the original proof reference \url{https://alpo.ge/remark.pdf}.
\bibitem{incoming} V. Reshetnikov, ProveIt incoming research reports.\newline \url{https://github.com/VladimirReshetnikov/ProveIt/tree/main/docs/incoming}.
\bibitem{heilbronn} H. Heilbronn, On real zeros of Dedekind $\zeta$-functions, Canad. J. Math. 25 (1973), 870--873.
\bibitem{foote} R. Foote, H. Ginsberg, V. K. Murty, Heilbronn characters, Bull. Amer. Math. Soc. 52 (2015), 465--496, \href{https://doi.org/10.1090/BULL/1492}{doi:10.1090/BULL/1492}.
\bibitem{carter} R. W. Carter, Conjugacy classes in the Weyl group, Compositio Math. 25 (1972), 1--59.
\bibitem{springer} T. A. Springer, Regular elements of finite reflection groups, Invent. Math. 25 (1974), 159--198, Theorem 8.7.
\bibitem{hurwitz} B. Baumeister, M. Dyer, C. Stump, P. Wegener, On the Hurwitz action in finite Coxeter groups, J. Group Theory 20 (2017), 103--131, \href{https://doi.org/10.1515/jgth-2016-0025}{doi:10.1515/jgth-2016-0025}. For an explicit length-additive formulation see \url{https://noah.nrw/ubbihs/download/pdf/5142443}, Lemma 4.2.10 and following discussion.
\bibitem{mv} H. L. Montgomery and R. C. Vaughan, \emph{Multiplicative Number Theory I: Classical Theory}, Cambridge, 2007, Chapter 15 (Landau positivity).
\bibitem{neukirch} J. Neukirch, \emph{Algebraic Number Theory}, Springer, 1999, chapters on zeta functions and class field theory.
\end{thebibliography}
\end{document}
